% Required Packages: \usepackage{graphicx}, \usepackage{float}

\begin{figure}[htbp]
  \centering
  \begin{subfigure}[b]{0.48\linewidth}
    \centering
    \zimg{rf_fig_8.png}{width=\linewidth,max height=0.4\textheight,keepaspectratio}{fg_0_r1w9m}{rf_fig_8.png}
    \caption{(b) shows that the output resistance in all TC-OTAs diminishes as N grows. Since an increase in N is equivalent to an increase...}

  \end{subfigure}
  \hfill
  \begin{subfigure}[b]{0.48\linewidth}
    \centering
    \zimg{rf_fig_9.png}{width=\linewidth,max height=0.4\textheight,keepaspectratio}{fg_1_imi8a}{rf_fig_9.png}
    \caption{(b) shows that the output resistance in all TC-OTAs diminishes as N grows. Since an increase in N is equivalent to an increase...}

  \end{subfigure}
  \caption{(b) shows that the output resistance in all TC-OTAs diminishes as N grows. Since an increase in N is equivalent to an increase in a transistor's width W, output resistance (ROUT) must drop as N increases. However, TC-OTAs have a greater ROUT due to the cascoding effect. Figure 3(c) shows how N affects the average power consumption of CNT-based TC-OTAs. It is evident that average power grows along with N. This is explained by the fact that CNTFETs' driving capacity has increased due to the increase in N. However, because to the exceptionally low CNT channel resistance, 1D ballistic transport, reduced parasitic capacitance, and reduced leakage, a pure CNT-based TC-OTA has a much lower power dissipation \cite{ref18,ref35}. When comparing the two hybrid TC-OTAs, the NCNTFET-PMOS-TC-OTA uses more power than the PCNTFET-NMOS-TC-OTA because the former has more nCNTFETs (14 nCNTFETs) than the latter (12 pCNTFETs). The bandwidth is greatly enhanced in CNT-based TC-OTAs, as shown in Figure 3(d), with the amount of CNTs having the biggest impact on pure CNT-TC-OTA. It is more important for PCNTFET-NMOS-TC-OTA circuits than NCNTFET-PMOS-TC-OTA circuits, though, and grows linearly \cite{ref6,ref19}.}

  \label{fig:group_67bqx}
\end{figure}
